%%=============================================================================
%% CAPITULO 1 -- INTRODUCAO
%% Esqueleto provisorio (dissertacao dummy): texto a ser desenvolvido.
%%=============================================================================

\section{Introdução}

O Canal Educação, programa de mediação tecnológica da Secretaria de Estado da
Educação do Piauí (SEDUC-PI), acumulou em mais de uma década de operação
milhares de videoaulas que cobrem etapas e modalidades da educação básica nos
224 municípios do estado. Dentro desse acervo, o conjunto de 760 videoaulas
ativas do Preparatório ENEM, distribuído em quinze disciplinas, constitui um
ativo digital público de escala considerável, cujo valor depende menos de sua
existência do que da capacidade institucional de convertê-lo em uso pedagógico
efetivo.

O professor da rede pública que identifica a necessidade de trabalhar
determinada habilidade da matriz de referência do Exame Nacional do Ensino
Médio (ENEM) não dispõe de ferramenta que indique, de forma automática e
contextualizada, quais aulas desse acervo são mais adequadas para esse fim. O
diagnóstico avaliativo e o acervo de conteúdo foram construídos por trilhas
institucionais distintas e ainda não possuem interfaces que os conectem no
momento do planejamento, de modo que a tarefa de articular currículo, avaliação
e acervo recai integralmente sobre o docente.

Esta pesquisa desenvolve e avalia, sob o método da Design Science Research
(DSR), um artefato de recomendação pedagógica que classifica as videoaulas do
Canal Educação segundo as 120 habilidades da matriz de referência do ENEM. O
artefato converte o conteúdo audiovisual em texto por transcrição automática,
representa aulas e descritores por embeddings do modelo multilíngue BGE-M3 e
ordena, para cada aula, as habilidades mais próximas pela similaridade de
cosseno, sem depender de histórico de uso.

\subsection{Objetivo geral}

Avaliar tecnicamente alternativas de classificação automatizada do acervo de
videoaulas do Preparatório ENEM do Canal Educação segundo as habilidades da
matriz de referência do ENEM, tomando como referência externa a atribuição
oficial de habilidades a itens do exame publicada pelo Instituto Nacional de
Estudos e Pesquisas Educacionais Anísio Teixeira (INEP).

\subsection{Objetivos específicos}

\begin{itemize}
  \item Mapear e classificar as 760 videoaulas ativas do Preparatório ENEM por
        habilidade da matriz de referência, caracterizando previamente acervo e
        matriz como ativos digitais do Estado;
  \item Construir o artefato de recomendação baseado em conteúdo, com
        embeddings BGE-M3 e similaridade de cosseno, sem dependência de
        histórico de uso;
  \item Comparar, sobre os itens do ENEM com habilidade atribuída pelo INEP, o
        método do artefato, um método lexical sem inteligência artificial
        (TF-IDF) e um modelo generativo, quanto a qualidade, estabilidade e
        viabilidade de execução.
\end{itemize}

\subsection{Problema}

Em que medida alternativas de classificação automatizada, com e sem
inteligência artificial, conseguem atribuir às videoaulas de um acervo público
as habilidades da matriz de referência do ENEM com concordância suficiente, em
relação à atribuição oficial do INEP, para apoiar o planejamento docente sem
substituir a decisão pedagógica?

\subsection{Justificativa}

Do ponto de vista da gestão pública, a classificação do acervo converte-se em
insumo de governança do próprio ativo: ao expor as habilidades de alta
recorrência no exame com baixa ou nenhuma videoaula disponível, orienta a
priorização da produção de novos conteúdos pela SEDUC-PI. O recorte declara,
ainda, o conflito de interesse do pesquisador, que exerce função de direção na
organização que opera o Canal Educação, e adota como salvaguarda a natureza
externa da referência de avaliação, produzida pelo INEP e anterior à pesquisa.
